% Relativistic Chapter I: Basic Concepts
% Relativistic generalization of Chandrasekhar, Chapter I (pp.\ 1--8)
% Conventions: see RELATIVISTIC_CONVENTIONS.md

\chapter{Basic Concepts: Relativistic Generalization}
\label{ch:rel-1}

%% ----------------------------------------------------------------
\section{Introduction}
\label{sec:rel-1-1}

The equations of relativistic hydrodynamics---like their
non-relativistic counterparts---admit simple stationary flow
patterns whose realizability is limited to certain ranges of the
governing parameters.  Outside these ranges the flows are
\emph{unstable}: they cannot sustain themselves against small
perturbations.  The demarcation of stable from unstable states
constitutes the subject of \emph{relativistic hydrodynamic stability}.

When the characteristic velocities in the problem approach the speed
of light~$c$, or when thermal and internal energies become comparable
to the rest-mass energy, the Navier--Stokes description of viscous
flow must be replaced by a manifestly covariant framework.  Two
physical requirements distinguish the relativistic theory from the
classical one:
\begin{enumerate}
\item \textbf{Covariance.}  All equations are tensor equations on
  spacetime, written in terms of the energy-momentum tensor
  $T^{\mu\nu}$, the baryon-number current $n\, u^{\mu}$, and
  (when electromagnetic fields are present) the Faraday tensor
  $F^{\mu\nu}$.
\item \textbf{Causality.}  No signal---including perturbation modes,
  thermal waves, and viscous diffusion fronts---may propagate faster
  than~$c$.  This requirement rules out parabolic transport equations
  (Fourier, Navier--Stokes) in favour of hyperbolic, causal
  alternatives such as the Israel--Stewart theory.
\end{enumerate}

In what follows we reformulate the basic concepts of Chapter~I in
this relativistic setting.  Throughout, Greek indices
$\mu,\nu,\ldots$ run over $0,1,2,3$ with signature $(-,+,+,+)$;
Latin indices $i,j,\ldots$ label spatial directions $1,2,3$.  The
speed of light~$c$ is kept explicit so that the non-relativistic limit
$v/c\to 0$ is transparent.  The four-velocity is normalised as
$u^{\mu}u_{\mu}=-c^{2}$, and the projection tensor orthogonal to the
flow is $\Delta^{\mu\nu}=g^{\mu\nu}+u^{\mu}u^{\nu}/c^{2}$.

%% ----------------------------------------------------------------
\section{Relativistic fluid equations}
\label{sec:rel-1-2}

\subsection{Perfect-fluid (Euler) sector}
\label{sec:rel-1-2a}

For a perfect fluid the energy-momentum tensor is
\begin{equation}
\label{eq:rel-1-1}
T^{\mu\nu} = \frac{\varepsilon + p}{c^{2}}\,u^{\mu}\,u^{\nu}
              + p\,g^{\mu\nu},
\end{equation}
where $\varepsilon$ is the total energy density (rest-mass plus
internal) and $p$ is the isotropic pressure.  The relativistic
enthalpy density is $w=(\varepsilon+p)/c^{2}$.  Conservation of
energy-momentum,
\begin{equation}
\label{eq:rel-1-2}
\nabla_{\mu}T^{\mu\nu}=0,
\end{equation}
together with baryon-number conservation,
\begin{equation}
\label{eq:rel-1-3}
\nabla_{\mu}(n\,u^{\mu})=0,
\end{equation}
yield the relativistic Euler equations.  Projecting
\eqref{eq:rel-1-2} along $u_{\nu}$ and orthogonal to $u_{\nu}$, one
obtains, respectively, the energy equation and the relativistic
momentum equation:
\begin{align}
u^{\mu}\nabla_{\mu}\varepsilon
  + (\varepsilon+p)\,\nabla_{\mu}u^{\mu} &= 0,
\label{eq:rel-1-4}\\[6pt]
w\,u^{\mu}\nabla_{\mu}u^{\nu}
  + \Delta^{\nu\alpha}\nabla_{\alpha}p &= 0.
\label{eq:rel-1-5}
\end{align}

\noindent\textbf{Non-relativistic limit.}  In the limit
$v/c\to 0$, $\varepsilon\to\rho c^{2}$, $w\to\rho$, the time-like
component of \eqref{eq:rel-1-5} reduces to the classical Euler
equation $\rho(\partial\boldsymbol{v}/\partial t +
\boldsymbol{v}\cdot\nabla\boldsymbol{v})=-\nabla p$, and the
conservation law~\eqref{eq:rel-1-3} becomes the continuity equation
$\partial\rho/\partial t+\nabla\cdot(\rho\boldsymbol{v})=0$.

\subsection{Israel--Stewart dissipative sector}
\label{sec:rel-1-2b}

In the presence of viscosity and heat conduction the energy-momentum
tensor acquires dissipative corrections:
\begin{equation}
\label{eq:rel-1-6}
T^{\mu\nu} = \varepsilon\,\frac{u^{\mu}u^{\nu}}{c^{2}}
  + (p+\Pi)\,\Delta^{\mu\nu}
  + \pi^{\mu\nu}
  + \frac{1}{c^{2}}(q^{\mu}u^{\nu}+q^{\nu}u^{\mu}),
\end{equation}
where $\Pi$ is the bulk viscous pressure, $\pi^{\mu\nu}$ is the
traceless shear stress tensor satisfying
$\pi^{\mu\nu}u_{\nu}=0$ and $\pi^{\mu}{}_{\mu}=0$, and $q^{\mu}$ is
the heat-flux four-vector with $q^{\mu}u_{\mu}=0$.

The Navier--Stokes (first-order) constitutive relations $\Pi=-\zeta\,
\nabla_{\mu}u^{\mu}$, $\pi^{\mu\nu}=2\eta_{\rm s}\,\sigma^{\mu\nu}$,
$q^{\mu}=-\kappa\,\Delta^{\mu\nu}\nabla_{\nu}T$ are \emph{acausal}:
they admit infinite propagation speeds for viscous and thermal
disturbances.  The Israel--Stewart (second-order) theory restores
causality by promoting each dissipative flux to an independent
dynamical variable obeying a relaxation equation:
\begin{align}
\tau_{\Pi}\,u^{\alpha}\nabla_{\alpha}\Pi + \Pi
  &= -\zeta\,\nabla_{\mu}u^{\mu},
\label{eq:rel-1-7}\\[4pt]
\tau_{\pi}\,u^{\alpha}\nabla_{\alpha}\pi^{\langle\mu\nu\rangle}
  + \pi^{\mu\nu}
  &= 2\eta_{\rm s}\,\sigma^{\mu\nu},
\label{eq:rel-1-8}\\[4pt]
\tau_{q}\,u^{\alpha}\nabla_{\alpha}q^{\langle\mu\rangle}
  + q^{\mu}
  &= -\kappa\!\left(
    \Delta^{\mu\nu}\nabla_{\nu}T
    + \frac{T}{c^{2}}\,u^{\mu}\,u^{\alpha}\nabla_{\alpha}T
  \right),
\label{eq:rel-1-9}
\end{align}
where $\tau_{\Pi}$, $\tau_{\pi}$, $\tau_{q}>0$ are the relaxation
times for bulk viscosity, shear viscosity, and heat flux respectively,
$\eta_{\rm s}$ is the shear viscosity, $\zeta$ the bulk viscosity,
$\kappa$ the thermal conductivity, and $\sigma^{\mu\nu}$ is the
symmetric trace-free shear tensor
\begin{equation}
\label{eq:rel-1-10}
\sigma^{\mu\nu} = \Delta^{\mu\alpha}\Delta^{\nu\beta}
  \!\left(\tfrac{1}{2}\nabla_{\alpha}u_{\beta}
        +\tfrac{1}{2}\nabla_{\beta}u_{\alpha}
        -\tfrac{1}{3}\Delta_{\alpha\beta}\,\nabla_{\gamma}u^{\gamma}
  \right).
\end{equation}

\noindent\textbf{Non-relativistic limit.}  Setting $\tau_{\Pi},
\tau_{\pi},\tau_{q}\to 0$ (instantaneous relaxation) and taking
$v/c\to 0$ recovers the Navier--Stokes constitutive relations and the
Fourier heat law.

\noindent\textbf{Causality enforcement.}  The relaxation times
$\tau_{\Pi},\tau_{\pi},\tau_{q}$ guarantee that the characteristic
speeds of the resulting hyperbolic system are bounded by~$c$.
In particular, the thermal signal speed
$v_{\rm th}=\sqrt{\kappa T/(\tau_{q}\,w\,c^{2})}$ and the viscous
signal speed $v_{\rm visc}=\sqrt{\eta_{\rm s}/(\tau_{\pi}\,w)}$
remain sub-luminal whenever the relaxation times satisfy the
Israel--Stewart thermodynamic inequalities.

%% ----------------------------------------------------------------
\section{Normal-mode analysis in the relativistic context}
\label{sec:rel-1-3}

\subsection{Background state and perturbations}

Consider a stationary relativistic flow characterised by a background
energy-momentum tensor $\bar T^{\mu\nu}$, four-velocity
$\bar u^{\mu}$, energy density $\bar\varepsilon$, pressure $\bar p$,
and---when dissipation is present---background values of $\Pi$,
$\pi^{\mu\nu}$, and $q^{\mu}$.  Let the system be characterised by
a set of parameters $X_1,\ldots,X_j$ (including geometric dimensions,
characteristic velocities, gravitational and electromagnetic field
strengths, and thermodynamic coefficients).

We perturb every dynamical field:
\begin{equation}
\label{eq:rel-1-11}
\varepsilon = \bar\varepsilon + \delta\varepsilon, \quad
p = \bar p + \delta p, \quad
u^{\mu} = \bar u^{\mu} + \delta u^{\mu}, \quad
\Pi = \bar\Pi + \delta\Pi, \quad
\pi^{\mu\nu} = \bar\pi^{\mu\nu} + \delta\pi^{\mu\nu}, \quad
q^{\mu} = \bar q^{\mu} + \delta q^{\mu},
\end{equation}
and the full energy-momentum perturbation is
\begin{equation}
\label{eq:rel-1-12}
\delta T^{\mu\nu}
  = \delta\!\left[
    \varepsilon\,\frac{u^{\mu}u^{\nu}}{c^{2}}
    + (p+\Pi)\,\Delta^{\mu\nu}
    + \pi^{\mu\nu}
    + \frac{1}{c^{2}}(q^{\mu}u^{\nu}+q^{\nu}u^{\mu})
  \right].
\end{equation}

The four-velocity perturbation $\delta u^{\mu}$ must respect the
normalisation constraint $u^{\mu}u_{\mu}=-c^{2}$, which to first
order requires
\begin{equation}
\label{eq:rel-1-13}
\bar u_{\mu}\,\delta u^{\mu} = 0,
\end{equation}
so that $\delta u^{\mu}$ has only three independent spatial
components in the local rest frame.

\subsection{Linearised equations and mode decomposition}

Substituting \eqref{eq:rel-1-11} into the conservation laws
\eqref{eq:rel-1-2}--\eqref{eq:rel-1-3} and the Israel--Stewart
equations \eqref{eq:rel-1-7}--\eqref{eq:rel-1-9}, and retaining
only first-order terms, yields a system of linearised equations for
$(\delta\varepsilon, \delta u^{\mu}, \delta\Pi, \delta\pi^{\mu\nu},
\delta q^{\mu})$.

Following the same logic as in the classical theory
(Chapter~\ref{ch:1}, \S\,3), we decompose an arbitrary perturbation
into normal modes.  For a system confined between two parallel planes
with background quantities depending only on the coordinate $z$
perpendicular to the planes:
\begin{equation}
\label{eq:rel-1-14}
\delta A(x^{\mu}) = \int_{-\infty}^{+\infty}\int_{-\infty}^{+\infty}
  dk_x\,dk_y\;
  \hat A_{\mathbf{k}}(z)\,
  \exp\!\left[i(k_x x + k_y y) + p_{\mathbf{k}}\,t\right],
\end{equation}
where $\delta A$ stands for any of the perturbed quantities,
$k=\sqrt{k_x^{2}+k_y^{2}}$ is the transverse wave number, and
$p_{\mathbf{k}}$ is the (complex) temporal eigenvalue.

For geometries with axial or spherical symmetry the mode decomposition
proceeds exactly as in the classical case---via azimuthal Fourier modes
$e^{im\theta}$ or spherical harmonics $Y_l^m(\theta,\varphi)$---but
in each case the mode amplitudes are now tuples involving
$(\delta\varepsilon, \delta u^{\mu}, \delta\Pi,
\delta\pi^{\mu\nu}, \delta q^{\mu})$ rather than the classical
velocity, temperature, and pressure perturbations alone.

\subsection{Eigenvalue problem and stability criterion}

The linearised perturbation equations, after mode decomposition,
reduce to an eigenvalue problem for the complex
frequency~$p_{\mathbf{k}}$:
\begin{equation}
\label{eq:rel-1-15}
p_{\mathbf{k}} = p_{\mathbf{k}}^{(r)} + i\,p_{\mathbf{k}}^{(i)}.
\end{equation}
The stability criterion is identical in form to the classical one:
\begin{itemize}
\item \textbf{Stable:} $p_{\mathbf{k}}^{(r)} < 0$ for every
  mode~$\mathbf{k}$.
\item \textbf{Unstable:} $\exists\;\mathbf{k}$ with
  $p_{\mathbf{k}}^{(r)}>0$.
\item \textbf{Marginal:}
  $p_{\mathbf{k}}^{(r)}(X_1,\ldots,X_j)=0$ defines the marginal
  locus $\Sigma_{\mathbf{k}}$ for mode~$\mathbf{k}$; the overall
  marginal surface~$\Sigma$ is the envelope of all
  $\Sigma_{\mathbf{k}}$.
\end{itemize}
The classification of marginal states into \emph{stationary}
($p_{\mathbf{k}}^{(i)}=0$, exchange of stabilities) and
\emph{oscillatory} ($p_{\mathbf{k}}^{(i)}\neq 0$, overstability)
carries over unchanged.

\subsection{Causality constraint on modes}
\label{sec:rel-1-3d}

A crucial new ingredient in the relativistic theory is the causality
bound on the phase and group velocities of every perturbation mode.
Defining the phase velocity
\begin{equation}
\label{eq:rel-1-16}
v_{\rm ph} = \frac{|p_{\mathbf{k}}^{(i)}|}{|\mathbf{k}|}
\end{equation}
and the group velocity
\begin{equation}
\label{eq:rel-1-17}
v_{\rm gr} = \left|
  \frac{\partial\,p_{\mathbf{k}}^{(i)}}{\partial \mathbf{k}}
\right|,
\end{equation}
the requirement of relativistic causality is
\begin{equation}
\label{eq:rel-1-18}
\boxed{v_{\rm ph}\leq c \quad\text{and}\quad v_{\rm gr}\leq c
  \qquad\text{for every mode } \mathbf{k}.}
\end{equation}
Any mode that violates \eqref{eq:rel-1-18} is physically inadmissible
and must be discarded.  In the Israel--Stewart framework
(\S\,\ref{sec:rel-1-2b}) this bound is \emph{automatically}
satisfied when the relaxation times are positive and the
thermodynamic stability conditions
$c_{s}^{2}=(\partial p/\partial\varepsilon)_{s}\geq 0$ and
$c_{s}^{2}\leq c^{2}$ hold.  This is the principal advantage of the
Israel--Stewart formulation over the acausal Navier--Stokes theory.

In non-dissipative systems the characteristic speeds are the sound
speed $c_s$, the Alfv\'en speed $v_A$, and the fast and slow
magnetosonic speeds $v_f$, $v_s$, all of which are bounded by~$c$
(see \S\,\ref{sec:rel-1-4}).

\noindent\textbf{Non-relativistic limit.}  As $v/c\to 0$ the
causality bound becomes trivially satisfied, and the eigenvalue
problem reduces to that of Chapter~\ref{ch:1}, \S\,3.

%% ----------------------------------------------------------------
\section{Relativistic dimensionless numbers}
\label{sec:rel-1-4}

The classical dimensionless numbers that govern hydrodynamic stability
acquire relativistic corrections when the flow velocities, sound
speeds, or thermal energies are no longer negligible compared to~$c$.
We define the relativistic generalisations below, with the Lorentz
factor $\gamma=(1-v^{2}/c^{2})^{-1/2}$ and the relativistic enthalpy
per unit volume $w=(\varepsilon+p)/c^{2}$.

\subsection{Reynolds number}

In the classical theory the Reynolds number is
$\mathrm{Re}=Lv/\nu$.  Relativistically, kinematic viscosity must be
replaced by the ratio $\nu_{\rm rel}=\eta_{\rm s}/(w\,c^{2})$
(where $\eta_{\rm s}$ is the dynamic shear viscosity), giving
\begin{equation}
\label{eq:rel-1-19}
\mathrm{Re}_{\rm rel}
  = \frac{L\,\gamma\,v}{\nu_{\rm rel}}
  = \frac{(\varepsilon+p)\,L\,\gamma\,v}{\eta_{\rm s}\,c^{2}}.
\end{equation}
When $v\ll c$ and $\varepsilon+p\approx\rho c^{2}$, this reduces to
$\mathrm{Re}=Lv/\nu$ with $\nu=\eta_{\rm s}/\rho$.

\subsection{Prandtl number}

The Prandtl number compares viscous to thermal diffusivity.
Relativistically:
\begin{equation}
\label{eq:rel-1-20}
\mathrm{Pr}_{\rm rel}
  = \frac{\eta_{\rm s}\,c_{p}}{\kappa}
  \left(1 + \frac{v^{2}}{c^{2}}\,\mathcal{O}(1)\right),
\end{equation}
where $c_{p}$ is the specific heat at constant pressure and $\kappa$
the thermal conductivity.  The leading-order expression is identical
to the classical Prandtl number; corrections enter at order
$v^{2}/c^{2}$ and are model-dependent.

\subsection{Rayleigh number}

For a horizontal layer of relativistic fluid of depth~$d$ heated from
below with adverse temperature gradient~$\beta$, the relativistic
Rayleigh number is
\begin{equation}
\label{eq:rel-1-21}
\mathrm{Ra}_{\rm rel}
  = \frac{g\,\alpha\,\beta\,d^{4}}{\kappa_{T}\,\nu_{\rm rel}}
  \left[
    1 + \mathcal{C}_{\rm Ra}\,\frac{v^{2}}{c^{2}}
    + \mathcal{O}\!\left(\frac{v^{4}}{c^{4}}\right)
  \right],
\end{equation}
where $\alpha$ is the relativistic expansion coefficient,
$\kappa_{T}$ the thermal diffusivity, $g$ the gravitational
acceleration, and $\mathcal{C}_{\rm Ra}$ is a dimensionless constant
depending on the equation of state.  In the non-relativistic limit
this reduces to the classical Rayleigh number
$R=g\alpha\beta d^{4}/(\kappa\nu)$ of equation~(14) of
Chapter~\ref{ch:1}.

\subsection{Taylor number}

For a rotating relativistic fluid the Taylor number generalises to
\begin{equation}
\label{eq:rel-1-22}
\mathrm{Ta}_{\rm rel}
  = \frac{4\,\Omega^{2}\,L^{4}}{\nu_{\rm rel}^{2}}
  \left[
    1 + \mathcal{C}_{\rm Ta}\,\frac{\Omega^{2}L^{2}}{c^{2}}
    + \mathcal{O}\!\left(\frac{v^{4}}{c^{4}}\right)
  \right],
\end{equation}
where the correction term $\mathcal{C}_{\rm Ta}$ encodes both
special-relativistic kinematic effects and (in general relativity)
frame-dragging contributions from the Lense--Thirring effect.

\subsection{Chandrasekhar number (magnetic field)}

When a magnetic field $b^{\mu}$ (the covariant magnetic four-vector in
the fluid frame) is present, the Chandrasekhar number generalises to
\begin{equation}
\label{eq:rel-1-23}
Q_{\rm rel}
  = \frac{b^{2}\,d^{2}}{4\pi\,w\,\nu_{\rm rel}\,\eta_m}
  \left[
    1 + \mathcal{O}\!\left(\frac{v_{A}^{2}}{c^{2}}\right)
  \right],
\end{equation}
where $b^{2}=b^{\mu}b_{\mu}$ is the magnetic pressure parameter,
$\eta_m$ is the magnetic diffusivity, and $v_{A}$ is the Alfv\'en
speed:
\begin{equation}
\label{eq:rel-1-24}
v_{A}^{2} = \frac{b^{2}}{4\pi\,w + b^{2}/c^{2}}\,c^{2}
  < c^{2}.
\end{equation}
The strict inequality in \eqref{eq:rel-1-24} is a consequence of the
causality bound \eqref{eq:rel-1-18}.

\subsection{Characteristic speeds and causality}
\label{sec:rel-1-4f}

All characteristic wave speeds in the relativistic theory are
sub-luminal by construction:
\begin{align}
\text{Sound:} &\quad c_{s}^{2}
  = \left(\frac{\partial p}{\partial\varepsilon}\right)_{\!s}
  \leq c^{2},
\label{eq:rel-1-25}\\[4pt]
\text{Alfv\'en:} &\quad v_{A}^{2}
  = \frac{b^{2}\,c^{2}}{4\pi\,w\,c^{2}+b^{2}} \leq c^{2},
\label{eq:rel-1-26}\\[4pt]
\text{Fast magnetosonic:} &\quad v_{f}^{2}
  = c_{s}^{2}+v_{A}^{2}-\frac{c_{s}^{2}\,v_{A}^{2}}{c^{2}}
  \leq c^{2},
\label{eq:rel-1-27}\\[4pt]
\text{Slow magnetosonic:} &\quad v_{s}^{2}
  = \frac{c_{s}^{2}\,v_{A}^{2}}{v_{f}^{2}}
  \leq c^{2}.
\label{eq:rel-1-28}
\end{align}
The fast magnetosonic bound $v_f^2\leq c^2$ follows from the
algebraic identity
$v_f^2 = c_s^2+v_A^2(1-c_s^2/c^2) \leq c_s^2+(c^2-c_s^2)=c^2$.

\subsection{Mach number}
The relativistic Mach number
\begin{equation}
\label{eq:rel-1-29}
\mathcal{M} = \frac{\gamma\,v}{c_{s}/\sqrt{1-c_{s}^{2}/c^{2}}}
\end{equation}
governs the importance of compressibility.  For $c_{s}\ll c$ and
$v\ll c$ this reduces to the classical $\mathcal{M}=v/c_{s}$.

%% ----------------------------------------------------------------
\section{Summary: recovery of the non-relativistic limit}
\label{sec:rel-1-5}

Every result of this chapter reduces to the corresponding result of
the classical Chapter~I when the following limits are taken
simultaneously:
\begin{enumerate}
\item $v/c\to 0$ and $\gamma\to 1$ (non-relativistic flow);
\item $\varepsilon\to\rho c^{2}$, $w\to\rho$ (rest-mass dominated
  energy);
\item $c_{s}/c\to 0$ (non-relativistic sound speed);
\item $\tau_{\Pi},\tau_{\pi},\tau_{q}\to 0$ (instantaneous
  relaxation, recovering Navier--Stokes and Fourier);
\item $v_{A}/c\to 0$ (non-relativistic Alfv\'en speed).
\end{enumerate}

Under these limits:
\begin{itemize}
\item The relativistic Euler equations
  \eqref{eq:rel-1-4}--\eqref{eq:rel-1-5} $\to$ classical Euler
  equations.
\item The Israel--Stewart equations
  \eqref{eq:rel-1-7}--\eqref{eq:rel-1-9} $\to$ Navier--Stokes +
  Fourier constitutive relations.
\item The normal-mode eigenvalue problem (\S\,\ref{sec:rel-1-3}) $\to$
  the eigenvalue problem of Chapter~\ref{ch:1}, \S\,3.
\item The causality bound \eqref{eq:rel-1-18} becomes vacuous.
\item $\mathrm{Re}_{\rm rel}\to\mathrm{Re}$,
  $\mathrm{Pr}_{\rm rel}\to\mathrm{Pr}$,
  $\mathrm{Ra}_{\rm rel}\to\mathrm{Ra}$,
  $\mathrm{Ta}_{\rm rel}\to\mathrm{Ta}$,
  $Q_{\rm rel}\to Q$.
\end{itemize}

%% ----------------------------------------------------------------
\section*{Bibliographical Notes}
\label{sec:rel-1-bib}

The Israel--Stewart theory of causal dissipative relativistic fluids
is developed in:

\begin{enumerate}
\item W.~\textsc{Israel}, `Nonstationary irreversible thermodynamics:
  a causal relativistic theory', \emph{Ann.\ Phys.\ (N.Y.)}
  \textbf{100}, 310--331 (1976).

\item W.~\textsc{Israel} and J.~M.~\textsc{Stewart}, `Transient
  relativistic thermodynamics and kinetic theory',
  \emph{Ann.\ Phys.\ (N.Y.)} \textbf{118}, 341--372 (1979).
\end{enumerate}

\noindent
For relativistic hydrodynamic stability and characteristic speeds:

\begin{enumerate}\setcounter{enumi}{2}
\item J.~L.~\textsc{Synge}, \emph{The Relativistic Gas}, North-Holland,
  Amsterdam, 1957.

\item A.~M.~\textsc{Anile}, \emph{Relativistic Fluids and
  Magneto-fluids}, Cambridge University Press, Cambridge, 1989.

\item P.~\textsc{Romatschke} and U.~\textsc{Romatschke},
  \emph{Relativistic Fluid Dynamics In and Out of Equilibrium},
  Cambridge University Press, Cambridge, 2019.
\end{enumerate}

\noindent
The relativistic MHD characteristic speed structure is discussed in:

\begin{enumerate}\setcounter{enumi}{5}
\item A.~\textsc{Lichnerowicz}, \emph{Relativistic Hydrodynamics and
  Magnetohydrodynamics}, W.~A.~Benjamin, New York, 1967.

\item S.~S.~\textsc{Komissarov}, `On the properties of Alfv\'en
  waves in relativistic magnetohydrodynamics',
  \emph{Phys.\ Lett.\ A} \textbf{232}, 435--442 (1997).
\end{enumerate}
